Badminton is a fast-paced racket sport in which rally dynamics can shift within
a handful of exchanges. Coaches and analysts regularly pose questions such as
which strokes dominate a match, what replies tend to follow a smash, whether
forward movement correlates with attacking opportunities, and whether opposing
players exhibit distinct tactical profiles. These are data-mining questions, yet
broadcast video does not supply the structured event tables needed to answer
them. The raw signal must be converted from pixels into an ordered stroke
sequence, and that conversion requires a chain of perception models whose errors
accumulate before any mining step begins.

This conversion is technically demanding. The shuttlecock is small and often
motion-blurred; broadcast footage interleaves replays, close-ups, scoreboard
inserts, and non-rally segments that must be filtered out. The far-court player
appears at low resolution and may be partially occluded, degrading pose estimates
at exactly the positions where attacking strokes are most frequent. Stroke
classification is further complicated by the fine-grained taxonomy of badminton
labels: a reliable system must distinguish, for example, a net shot from a return
net, a drop from a push, or a smash from closely related attacking strokes.
Without accounting for this error propagation, tactical analyses risk drawing
conclusions from a corrupted event stream.

We construct and evaluate \method, a practical end-to-end pipeline that converts
badminton broadcast video into mineable tactical data. The overall architecture
is depicted in Figure~\ref{fig:workflow}. The central research question is:
\emph{\textbf{Can an integrated raw-video pipeline recover badminton stroke sequences
that are sufficiently accurate and continuous to support meaningful tactical
mining on unseen matches?}}  Our answer is cautiously affirmative: by combining
simple yet reliable perception modules with careful integration, we recover
coherent stroke sequences that support a variety of interpretable analyses.

It is important to state the scope of this benchmark. The pipeline processes
unseen raw frames through rally filtering, shuttle tracking, pose extraction,
feature construction, stroke classification, and tactical mining. It does not
yet constitute a fully automatic system for arbitrary unlabeled broadcasts: the
unseen-video benchmark relies on ShuttleSet hit-frame indices to construct
stroke windows and on ShuttleSet homographies for court normalization. We make
this boundary explicit because it defines what the reported metrics demonstrate,
and because the primary bottleneck we identify—far-side player extraction—remains
an open problem even under fully automatic setups.

Our main contributions are:
\begin{itemize}[leftmargin=*,nosep]
  \item an integrated raw-video badminton pipeline linking rally filtering,
  shuttle tracking, pose extraction, and stroke classification into a single
  reproducible workflow;
  \item five tactical-mining analyses—stroke composition, Markov transitions,
  frequent motifs, movement profiling, and clustering—that illustrate how the
  recovered event stream can support diverse insights;
  \item a mining-validation protocol that re-runs every analysis on matched
  ground-truth labels and reports the predicted-vs-true gap for each view
  (distribution KL/$\chi^2$, motif rank correlation, transition baselines, and
  cluster ARI), turning unverified tactical claims into measured ones.
\end{itemize}

\begin{figure*}[t]
  \centering
  \includegraphics[width=\textwidth]{images/Overall_workflow.png}
  \caption{\textbf{End-to-end \method{} pipeline.} Raw broadcast frames are
  processed through rally filtering, shuttle tracking, player pose extraction,
  and stroke classification. The resulting structured stroke table drives all
  five tactical mining analyses. In our benchmark, ShuttleSet hit-frame labels
  and homographies are supplied for window construction and court normalization.}
  \label{fig:workflow}
\end{figure*}
